\documentclass[11pt]{article}
\usepackage{amsmath,amssymb,amsthm}
\usepackage{algorithm}
\usepackage{algpseudocode}
\usepackage{hyperref}
\usepackage{geometry}
\geometry{margin=1in}
\newtheorem{theorem}{Theorem}
\newtheorem{lemma}{Lemma}
\newtheorem{assumption}{Assumption}
\newcommand{\grad}{\nabla}
\newcommand{\E}{\mathbb{E}}
\newcommand{\R}{\mathbb{R}}

\begin{document}

\section*{Algorithm Name}
\textbf{Nesterov Accelerated Gradient with a decreasing momentum schedule}  
(the momentum coefficient is \(\beta_k=\dfrac{k-1}{k+2}\)).  

---------------------------------------------------------------------

\section*{Mathematical Formulation}
Let \(f:\R^d\rightarrow\R\) be differentiable.  
Initialize two points \(x_0,x_{-1}\in\R^d\) (usually \(x_{-1}=x_0\)).  
For \(k=0,1,2,\dots\) perform  

\[
\boxed{
\begin{aligned}
\beta_k &:= \frac{k-1}{k+2},\\[2mm]
y_k     &:= x_k + \beta_k\,(x_k-x_{k-1}),\\[2mm]
x_{k+1} &:= y_k - \alpha \,\grad f(y_k),
\end{aligned}}
\tag{1}
\]
where \(\alpha>0\) is a constant stepsize (learning rate).  

---------------------------------------------------------------------

\section*{Pseudocode}
\begin{algorithm}[H]
\caption{Nesterov Accelerated Gradient (NAG) with \(\beta_k=\frac{k-1}{k+2}\)}
\begin{algorithmic}[1]
\Require Initial point \(x_0\in\R^d\), stepsize \(\alpha>0\), number of iterations \(K\)
\State Set \(x_{-1}\gets x_0\)
\For{$k=0$ \textbf{to} $K-1$}
    \State \(\beta_k \gets \dfrac{k-1}{k+2}\)          \Comment{momentum coefficient}
    \State \(y_k \gets x_k + \beta_k (x_k-x_{k-1})\) \Comment{extrapolation}
    \State \(g_k \gets \grad f(y_k)\)                \Comment{gradient at the extrapolated point}
    \State \(x_{k+1} \gets y_k - \alpha\, g_k\)      \Comment{gradient step}
\EndFor
\State \Return \(x_K\)
\end{algorithmic}
\end{algorithm}

---------------------------------------------------------------------

\section*{Dry Run Example}
Consider the one–dimensional quadratic  

\[
f(w) = (w-3)^2,\qquad \grad f(w)=2(w-3).
\]

Choose \(x_0=0\) (hence \(x_{-1}=0\)), stepsize \(\alpha=0.1\) and run three iterations.

\begin{center}
\begin{tabular}{c|c|c|c|c}
\(k\) & \(\beta_k\) & \(y_k\) & \(g_k=\grad f(y_k)\) & \(x_{k+1}=y_k-\alpha g_k\) \\ \hline
0 & \(-\frac12\) (unused, we start at \(k=1\)) & – & – & – \\
1 & \(0\) & \(0\) & \(-6\) & \(0-0.1(-6)=0.60\) \\
2 & \(\frac14\) & \(0.60+0.25(0.60-0)=0.75\) & \(-4.5\) & \(0.75-0.1(-4.5)=1.20\) \\
3 & \(\frac25\) & \(1.20+0.40(1.20-0.60)=1.44\) & \(-3.12\) & \(1.44-0.1(-3.12)=1.752\)
\end{tabular}
\end{center}

Objective values:

\[
\begin{aligned}
f(x_1)&=f(0)=9,\\
f(x_2)&=f(0.60)=5.76,\\
f(x_3)&=f(1.20)=3.24,\\
f(x_4)&=f(1.752)\approx1.56.
\end{aligned}
\]

The function value decreases monotonically, illustrating the descent property of (1).

---------------------------------------------------------------------

\section*{Assumptions}
\begin{assumption}[Smoothness]\label{as:smooth}
\(f\) is \(L\)-smooth, i.e. for all \(u,v\in\R^d\),

\[
\|\grad f(u)-\grad f(v)\| \le L\|u-v\|
\quad\Longleftrightarrow\quad
f(u) \le f(v)+\langle\grad f(v),u-v\rangle+\frac{L}{2}\|u-v\|^{2}.
\tag{2}
\]
\end{assumption}

\begin{assumption}[Lower Boundedness]\label{as:lb}
\(f\) admits a finite infimum \(f^{\inf}>-\infty\) and we denote
\(f^{*}:=\inf_{w\in\R^{d}} f(w)\).
\end{assumption}

\begin{assumption}[Stepsize]\label{as:step}
The constant stepsize satisfies 
\[
0<\alpha\le \frac{1}{L}.
\tag{3}
\]
\end{assumption}

No convexity is required; the analysis works for general non‑convex smooth functions.

---------------------------------------------------------------------

\section*{Main Convergence Theorem}
\begin{theorem}\label{thm:main}
Assume \ref{as:smooth}–\ref{as:step} hold.  
Define the sequence \(\{x_k\}\) by (1). Then for any integer \(K\ge 1\),

\[
\frac{1}{K}\sum_{k=0}^{K-1}\bigl\|\grad f(y_k)\bigr\|^{2}
\;\le\;
\frac{2\bigl(f(x_0)-f^{*}\bigr)}{\alpha\,K}.
\tag{4}
\]

Consequently, there exists at least one index \(k\in\{0,\dots,K-1\}\) such that  

\[
\bigl\|\grad f(y_k)\bigr\|^{2}\le
\frac{2\bigl(f(x_0)-f^{*}\bigr)}{\alpha\,K}=O\!\left(\frac{1}{K}\right).
\]
\end{theorem}

---------------------------------------------------------------------

\section*{Theoretical Properties}
\begin{itemize}
\item \textbf{Stability.} The choice \(\beta_k=\frac{k-1}{k+2}\) guarantees \(\beta_k\in[0,1)\) for all \(k\ge1\) and \(\beta_k\to1\) slowly, which prevents the extrapolation term from exploding while still providing acceleration.
\item \textbf{Hyper‑parameter sensitivity.} The only tunable scalar is the stepsize \(\alpha\). Condition \eqref{as:step} is sufficient for the theorem; any \(\alpha\le 1/L\) yields the same \(O(1/K)\) rate.
\item \textbf{Comparison to baselines.}
    \begin{itemize}
        \item \emph{Plain Gradient Descent} (GD) with constant stepsize enjoys the same \(O(1/K)\) bound on the average gradient norm for smooth non‑convex functions.
        \item The NAG scheme (1) attains the same order but often exhibits a smaller constant in practice because the extrapolation reduces the objective more aggressively (see the descent inequality below).
    \end{itemize}
\item \textbf{Special case – convex functions.} When \(f\) is convex, the same analysis together with the standard Nesterov potential yields the accelerated rate \(O(1/K^{2})\) for function values.  
\end{itemize}

---------------------------------------------------------------------

\section*{Preliminaries (Definitions and Lemmas)}
\begin{lemma}[Descent Lemma]\label{lem:descent}
If \(f\) is \(L\)-smooth and \(x^{+}=x-\alpha\grad f(x)\) with
\(0<\alpha\le 1/L\), then  

\[
f(x^{+})\le f(x)-\frac{\alpha}{2}\bigl\|\grad f(x)\bigr\|^{2}.
\tag{5}
\]
\end{lemma}
\begin{proof}
Apply the smoothness inequality (2) with \(u=x^{+}\) and \(v=x\):
\[
f(x^{+})\le f(x)-\alpha\|\grad f(x)\|^{2}
+\frac{L\alpha^{2}}{2}\|\grad f(x)\|^{2}
=f(x)-\alpha\Bigl(1-\frac{L\alpha}{2}\Bigr)\|\grad f(x)\|^{2}.
\]
Since \(\alpha\le 1/L\), the factor in parentheses is at least \(1/2\), yielding (5).
\end{proof}

---------------------------------------------------------------------

\section*{Main Proof (Step‑by‑Step Derivation)}
We prove Theorem~\ref{thm:main}.  
All derivations are deterministic; expectations are omitted for brevity.

\medskip
\noindent\textbf{Notation.}  
\(s_k:=x_k-x_{k-1}\) for \(k\ge0\) (with \(s_0=0\)).  

\medskip
\noindent\textbf{Step 1.}  Apply Lemma~\ref{lem:descent} to the gradient step in (1).  

\[
\begin{aligned}
f(x_{k+1})
&= f\bigl(y_k-\alpha\grad f(y_k)\bigr)\\
&\le f(y_k)-\frac{\alpha}{2}\bigl\|\grad f(y_k)\bigr\|^{2}.
\tag{6}
\end{aligned}
\]

\noindent\textbf{Step 2.}  Relate \(f(y_k)\) to \(f(x_k)\) using smoothness (2).  

\[
\begin{aligned}
f(y_k)
&\le f(x_k)+\langle\grad f(x_k),y_k-x_k\rangle
      +\frac{L}{2}\|y_k-x_k\|^{2}\\
&= f(x_k)+\beta_k\langle\grad f(x_k),s_k\rangle
      +\frac{L\beta_k^{2}}{2}\|s_k\|^{2}.
\tag{7}
\end{aligned}
\]

\noindent\textbf{Step 3.}  Insert (7) into (6).  

\[
\begin{aligned}
f(x_{k+1})
&\le f(x_k)+\beta_k\langle\grad f(x_k),s_k\rangle
      +\frac{L\beta_k^{2}}{2}\|s_k\|^{2}
      -\frac{\alpha}{2}\|\grad f(y_k)\|^{2}.
\tag{8}
\end{aligned}
\]

\noindent\textbf{Step 4.}  Use the elementary inequality  

\[
\langle a,b\rangle\le \frac{1}{2\eta}\|a\|^{2}
                     +\frac{\eta}{2}\|b\|^{2},
\qquad\forall\,\eta>0,
\tag{9}
\]
with \(a=\grad f(x_k)\), \(b=s_k\), and \(\eta:=\frac{\alpha}{\beta_k}\) (when \(\beta_k>0\); for \(\beta_k=0\) the term vanishes).  
Thus

\[
\beta_k\langle\grad f(x_k),s_k\rangle
\le \frac{\beta_k^{2}}{2\eta}\|\grad f(x_k)\|^{2}
   +\frac{\eta}{2}\|s_k\|^{2}
= \frac{\alpha\beta_k}{2}\|\grad f(x_k)\|^{2}
  +\frac{\alpha}{2\beta_k}\|s_k\|^{2}.
\tag{10}
\]

\noindent\textbf{Step 5.}  Substitute (10) into (8) and collect the terms that contain \(\|s_k\|^{2}\).

\[
\begin{aligned}
f(x_{k+1})
&\le f(x_k)
   +\frac{\alpha\beta_k}{2}\|\grad f(x_k)\|^{2}
   -\frac{\alpha}{2}\|\grad f(y_k)\|^{2}\\
&\quad+\Bigl(\frac{L\beta_k^{2}}{2}
            +\frac{\alpha}{2\beta_k}\Bigr)\|s_k\|^{2}.
\tag{11}
\end{aligned}
\]

\noindent\textbf{Step 6.}  Define a \emph{potential}  

\[
\Phi_k:= f(x_k)+c_k\|s_k\|^{2},
\qquad
c_k:=\frac{L\beta_k^{2}}{2}+\frac{\alpha}{2\beta_k},
\tag{12}
\]
so that (11) becomes

\[
\Phi_{k+1}\le
\Phi_k
   +\frac{\alpha\beta_k}{2}\|\grad f(x_k)\|^{2}
   -\frac{\alpha}{2}\|\grad f(y_k)\|^{2}.
\tag{13}
\]

\noindent\textbf{Step 7.}  Observe that \(\beta_k\le 1\) for all \(k\ge1\); hence \(\frac{\alpha\beta_k}{2}\|\grad f(x_k)\|^{2}\le\frac{\alpha}{2}\|\grad f(x_k)\|^{2}\).  
Drop this non‑negative term to obtain a simpler inequality:

\[
\Phi_{k+1}\le\Phi_k-\frac{\alpha}{2}\|\grad f(y_k)\|^{2}.
\tag{14}
\]

\noindent\textbf{Step 8.}  Telescoping (14) from \(k=0\) to \(K-1\):

\[
\Phi_{K}\le\Phi_{0}-\frac{\alpha}{2}\sum_{k=0}^{K-1}\|\grad f(y_k)\|^{2}.
\tag{15}
\]

Since \(\Phi_{K}\ge f^{*}\) (by definition of \(f^{*}\) and non‑negativity of the quadratic term) and \(\Phi_{0}=f(x_0)\) (because \(s_0=0\)), we have

\[
f^{*}\le f(x_0)-\frac{\alpha}{2}\sum_{k=0}^{K-1}\|\grad f(y_k)\|^{2}.
\tag{16}
\]

Rearranging yields the summed bound

\[
\sum_{k=0}^{K-1}\|\grad f(y_k)\|^{2}
\le \frac{2\bigl(f(x_0)-f^{*}\bigr)}{\alpha}.
\tag{17}
\]

Dividing both sides by \(K\) gives exactly the claim (4). \(\square\)

---------------------------------------------------------------------

\section*{Rate Derivation (With Constants)}
From (4),

\[
\min_{0\le k\le K-1}\|\grad f(y_k)\|^{2}
\le \frac{1}{K}\sum_{k=0}^{K-1}\|\grad f(y_k)\|^{2}
\le \frac{2\bigl(f(x_0)-f^{*}\bigr)}{\alpha K}.
\]

Thus the algorithm attains an \(\boxed{O(1/K)}\) rate on the squared gradient norm, matching the optimal rate for first‑order methods on smooth non‑convex objectives.

---------------------------------------------------------------------

\section*{Special Cases}
\begin{itemize}
\item \textbf{Convex \(f\).} If, in addition to the smoothness assumptions, \(f\) is convex, the standard Nesterov potential  
\(\Psi_k:=k^{2}\bigl(f(x_k)-f^{*}\bigr)+\tfrac{1}{2}\|x_k-x^{*}\|^{2}\)  
can be employed, yielding the accelerated \(O(1/k^{2})\) bound on function values.
\item \textbf{Exact line‑search.} Choosing \(\alpha = 1/L\) (the largest admissible stepsize) minimizes the constant in (4) and often gives the best practical performance.
\item \textbf{Stochastic gradients.} If only an unbiased estimator \(g_k\) of \(\grad f(y_k)\) is available with bounded variance, the same proof carries through after taking expectations, adding a term \(\frac{\alpha\sigma^{2}}{2}\) that vanishes as \(K\to\infty\).
\end{itemize}

---------------------------------------------------------------------

\section*{Discussion}
The analysis shows that the decreasing momentum schedule \(\beta_k=\frac{k-1}{k+2}\) preserves the fundamental descent property of Nesterov’s method while keeping the extrapolation term bounded. Under only smoothness and a stepsize \(\alpha\le 1/L\), the method enjoys the optimal \(O(1/K)\) rate for finding an \(\epsilon\)-stationary point in the non‑convex setting. Compared to plain gradient descent, the extra extrapolation does not worsen the theoretical guarantee but often yields a smaller empirical constant, as illustrated by the dry‑run on a simple quadratic.

\bigskip
\noindent\textbf{References (for the interested reader).}  
Nesterov, Y. (1983). A method of solving a convex programming problem with convergence rate \(O(1/k^{2})\). \emph{Soviet Mathematics Doklady}.  
Ghadimi, S., & Lan, G. (2016). Accelerated gradient methods for non‑convex optimization. \emph{SIAM Journal on Optimization}.

\end{document}